\section{Introduction}\label{sec:intro}
Since the Scotland Acts of 2012 and 2016 devolved and partially devolved taxes now fund close to half of the Scottish Government's resource spending, with non-savings, non-dividend (NSND) income tax the largest part \citep{sfc2026}. Under the block grant adjustment (BGA), now the block grant is cut by foregone revenue, and that deduction indexed thereafter to the growth of the equivalent rUK tax base \citep{gov2023}. The Scottish budget therefore gains only insofar as Scottish income tax grows faster than its rUK comparator, such that the level of Scottish revenue is immaterial and only the differential matters --- a pattern that underlies much of the theory and application of this paper. Forecast at £969 million for 2026--27, the Scottish Fiscal Commission (SFC) records this gap between devolved revenue and the block grant as the income tax net position. Thus, a relative deterioration in Scottish earnings shrinks the net position and since outturn is reconciled against forecast three years later, an error made now goes on constraining spending long afterwards \citep{gov2023}.

Artificial intelligence is a shock capable of impacting such a difference, both in the short and the long run. Any technology that raises productivity and alters labour demand unevenly across occupations will then also reshape earnings unevenly across the economy, with any resulting divergence between the Scottish and rUK tax bases surfacing in the aforementioned net position \citep{sfc2016}. There is good reason to expect such a divergence due to Scotland's distinctive occupational composition. The direction may culminate from the composition harbouring Scotland from displacement, or it may deny Scotland the gains that accrue to complemented work. This is furthered by a productivity gap that predates AI, with growth weak across Scotland, the UK, and much of Europe since the 2008 crisis meaning a worsened fiscal outlook \citep{sfc2026}. Whether AI narrows or widens that gap is a question that depends on how its gains are distributed across regions.

%The standard input to that judgement is now an exposure score produced by a large language model, an approach \citet{eloundou2024} established and institutional forecasters have since adopted, among them the \citet{obr2025} for the UK productivity outlook. Recent evidence unsettles the premise that such scores are objective measurements. \citet{yin2026} score one task set with four LLMs and find average exposure differs by more than a factor of three, with the labour-market figures built on those scores varying roughly 2.4-fold. This divergence is systematic not random, and so averaging does not remove it. What, then, can a fiscal exercise like this one take from a number that changes with the model that produced it? This dissertation addresses that question, and answers it for this fiscal application. The solution rests on a coincidence that the BGA responds to the Scotland--rUK \textit{difference} not the Scottish \textit{level}, and the exposure index is dependable in differences but not in levels.

To assess this, the standard input is now an exposure score produced by a large language model, an approach \citet{eloundou2024} established and institutional forecasters have since adopted, among them the \citet{obr2025} for the UK productivity outlook and the ILO for a global index \citep{gmyrek2023}. Those scores are reported as point estimates with dependence on the scoring model merely conceded, yet recent evidence unsettles the premise that they are objective measurements at all. \citet{yin2026} score one task set with four LLMs and find average exposure differs by more than a factor of three, with the labour-market figures built on those scores varying roughly 2.4-fold. This divergence is systematic not random, and so averaging does not remove it. What, then, can a fiscal exercise like this one take from a number that changes with the model that produced it? This dissertation addresses that question, and answers it for this fiscal application. No instrument for AI exposure yet matches the commuting-zone design of \citet{acemoglu2020}, so what follows is measurement rather than causal identification. The solution rests on a coincidence: the BGA responds to the Scotland--rUK \textit{difference} not the Scottish \textit{level}, and the exposure index is dependable in differences but not in levels.

%A technology takes over tasks capital can perform and reinstates labour in new ones \citep{acemoglu2019}, and so a single exposure score pools both channels.
%An exposure score carries two channelstechnology takes over tasks capital can perform and reinstates labour in new ones \citep{acemoglu2019}, so one score carries both channels at once. 
%Score construction has moved from occupation-level bottlenecks \citep{frey2017} and ability crosswalks \citep{felten2021} to language models rating the individual task \citep{eloundou2024}, with applications extending into institutional practice \citep{tbi2024, obr2025, gmyrek2023}. Reporting convention has stayed as exposure arriving as a point estimate whose model-dependence is conceded rather than measured. 
%, and the spatial work inherits that habit along with the number \citep{tbi2024, oecd2024}. No instrument yet matches the commuting-zone design of \citet{acemoglu2020}, so no causal design is available for this study. What remains available is measurement as the above studies do, whilst quantifying scoring noise and establishing where in the chain from score to downstream exercises it is and is not troublesome.

%The problem has a familiar econometric look. Exposure enters downstream estimates as a regressor measured with error, where noise attenuates and systematic bias contaminates \citep{bound1991}. The regional index is a shift-share object, so employment-share weights govern how any scoring error aggregates. Lastly, the central result works by differencing where the Scotland--rUK differential removes the fixed model bias just as a difference-in-differences removes a unit fixed effect. REMOVE? PLUS ABOVE?

Three contributions follow, with the first being analytical. Proposition 1 shows that the regional differential in an exposure index removes the common component of any model-specific bias, retaining only the part co-varying with the cross-region employment-share gap, with Corollary 1 bounding that part by a Cauchy--Schwarz inequality. Corollary 2 covers the case where a model disagrees about the scale of the index while agreeing about order, and there a rank-based differential removes the bias in full. I evaluate whether the result holds by re-scoring the full task set on two further language models. Secondly, I contribute a compositional answer to the fiscal question, locating Scotland fifth of the twelve UK regions and nations, attributing about two-fifths of the headline gap to London. Thirdly, I quantify the impact of AI in fiscal units, carrying scoring uncertainty through. 


Scotland's employment-weighted continuous exposure sits 0.86 index points below the rUK on the 0--100 scale, with a 95\% interval of $[-1.07, -0.69]$, with most of the width coming from survey sampling. The scale-free reading translates to about two percentile points of the UK occupational exposure distribution. The sign survives two model re-scores, the raw gap moving to $-0.53$ under Claude Haiku and $-0.43$ under GPT-4o as absolute levels shift by up to 9.1 index points.
%; most of that movement is the scale compression Corollary 2 describes, and the non-affine residual is 23--24\% of the Cauchy--Schwarz worst case. 
%The classified exposed share, the object closest to the policy literature's headline numbers, is censored at 91\% of employment and its sign is unstable to its own thresholds, so it cannot carry the regional question. 
Passed through the Acemoglu calibration, the differential implies a ten-year TFP differential of around three basis points, which \cref{sec:fiscal} translates into a net-position sensitivity of about \pounds\fiscCen~million a year. Scoring noise moves this figure by \fiscNoisePct\%, survey sampling and the adoption scenario by around a two-fifths each, and the choice of scoring model by \fiscModelPct\%.

The rest of the paper is organised as follows. \Cref{sec:background} sets out the fiscal backdrop, the task-based mechanism, and the measurement literature, ultimately proposing the three research gaps I address. \Cref{sec:data} describes the data. \Cref{sec:method} constructs the exposure scores, states the region-invariance assumption, proves the cancellation results, and sets out the calibration and propagation design. \Cref{sec:results} presents the differential, its anatomy, its robustness across thresholds and scoring models, and the exposure--wage gradient. \Cref{sec:translation} carries the differential into projected productivity and then into the income tax net position, and quantifies uncertainty in the resulting figure. \Cref{sec:conclusion} concludes. \cref{fig:procedure} depicts the paper's procedure from start to finish. 